\begin{abstract}
The Efficient Market Hypothesis has been revisited by Behavioral Finance, which highlights the role of informational flow and investors' bounded rationality in asset pricing. In emerging markets, highly liquid and politically exposed assets, such as Petrobras preferred shares (PETR4), exhibit pronounced sensitivity to news shocks. This dissertation investigates the impact of financial-news \textbf{sentiment}, defined as the polarity score in the $[-1,+1]$ interval assigned to a text by a classifier, on forecasting the direction and volatility of the PETR4 asset on the B3 exchange, from 2018 to 2025. An empirical-quantitative approach is adopted, structured as a reproducible pipeline. News items were collected directly from five outlets through their REST Application Programming Interface (API), which secured the \textbf{exact date and time} of each publication, an essential requirement for the causal \textit{Lead-Lag} alignment between the news and the trading session. The resulting corpus of 205,697 news items was organized by a supervised taxonomy of seven thematic categories, and its sentiment was extracted by \textbf{FinBERT-PT-BR}, a BERT-family model fine-tuned to the financial domain in Portuguese. Volatility was modeled by a GARCH(1,1) with Student's \textit{t} distribution, and direction forecasting employed Support Vector Machines (SVM) and Extreme Gradient Boosting (XGBoost) within a Data Fusion structure under a three-way chronological split. Results show that fusing sentiment with prices raises test accuracy from 50.9\% (baseline) to 54.5\%, a gain statistically above chance (binomial test, $p=0.024$). Granger causality confirms that sentiment leads returns by one day and, above all, leads volatility strongly and persistently. We conclude that sentiment contributes modestly but significantly to direction forecasting, and more clearly to volatility, consistent with the near-efficiency of the market.

\vspace{\onelineskip}
\noindent
\textbf{Keywords}: Behavioral Finance. Natural Language Processing. FinBERT-PT-BR. Machine Learning. GARCH. PETR4.
\end{abstract}
